\section{Results}
\label{sec:results}

\subsection{Approximation Ratio}

Across all 300~instances---3~datasets $\times$ 5~sizes $\times$ 20~seeds---the
greedy algorithm matched the exact Z3 solution on every instance.
The approximation ratio $\rho = w_{\text{greedy}} / w_{\text{exact}}$ was
exactly $1.00$ in all cases, with zero variance.
This result holds uniformly across $|C| \in \{10, 20, 50, 100, 200\}$ and is
not an artefact of solver failure: the exact success rate was $1.0$ throughout,
and brute-force enumeration independently confirmed Z3's optima on all 15
small instances ($|C| = 10$).

A ratio of $1.00$ on every instance does not imply greedy is generally optimal.
Rather, it reflects the structural simplicity of the three generator families,
as we analyse below.

\subsection{Structural Explanation by Dataset}

\paragraph{Dataset A.}
Each instance contains a single MUS of size~2 with equal-weight constraints.
Greedy discovers this MUS in one iteration, selects either element, and halts.
Because the MUS has only two elements of equal weight, any singleton correction
subset is optimal; greedy cannot diverge from exact.
Mean greedy iterations: $1.00$; mean MUSes discovered: $1.00$.

\paragraph{Dataset B.}
Each instance plants 4~constraints on a shared variable $x$ with weights
$1.0, 2.0, 3.0, 4.0$, producing $\binom{4}{2} = 6$ pairwise MUSes in the
planted core.
The lazy greedy run discovers 5~MUSes before the residual instance becomes
satisfiable.
In this generated clique-like structure, the greedy choices coincide with the
exact optimum.
Mean greedy iterations: $5.00$; mean MUSes discovered: $5.00$, constant across
sizes because the planted conflict structure does not grow with $|C|$.
Greedy runtime ($35.45$\,ms mean) exceeds exact runtime ($13.19$\,ms mean) on
this dataset because greedy performs 5~sequential MUS-extraction passes while
Z3 solves the optimisation encoding in a single call.

\paragraph{Dataset C.}
Each instance plants a constructor mismatch: $f$ is forced to be both a
function type and a list type.
The resulting MUS contains exactly the two conflicting constraints
(\texttt{C\_shape\_function} and \texttt{C\_conflict\_list}); the other planted
constraints ($\alpha = \texttt{Int}$, $\beta = \texttt{Bool}$) are consistent
with either resolution and do not participate in a minimal conflict.
Greedy discovers this MUS in one iteration and removes one of its elements.
Mean greedy iterations: $1.00$; mean MUSes discovered: $1.00$.

\subsection{Runtime Scaling}

Greedy runtime grows with $|C|$ across all datasets, driven by the cost of
unification checks during MUS extraction (Figure~\ref{fig:runtime}).
On Dataset~A, mean greedy runtime increases from $2.94$\,ms at $|C|=10$ to
$25.41$\,ms at $|C|=200$.
Dataset~C shows a similar profile ($1.41$\,ms to $30.61$\,ms).
Dataset~B is slower due to its 5-iteration structure, rising from $5.89$\,ms
to $92.64$\,ms.
Exact Z3 runtime is comparatively stable across sizes and datasets
($12$--$13$\,ms mean), suggesting that, on these small synthetic instances,
the fixed overhead of the Z3 optimisation call dominates the additional filler
constraints.

\begin{figure}[t]
  \centering
  \includegraphics[width=\linewidth]{figure1_runtime_scaling}
  \caption{Runtime scaling of greedy and exact Z3 baselines on log--log axes.
    Dataset~B is consistently slower for greedy due to its 5-iteration MUS
    discovery structure; Datasets~A and~C converge toward exact runtime at
    larger $|C|$.}
  \label{fig:runtime}
\end{figure}

Figure~\ref{fig:overhead} shows that greedy overhead relative to exact is
explained by MUS discovery count.
Datasets~A and~C discover exactly 1~MUS per instance regardless of $|C|$,
keeping the relative runtime at or below $1.0$ for larger sizes.
Dataset~B discovers exactly 5~MUSes per instance, one per greedy iteration,
and its relative runtime grows from about $2.0\times$ at $|C|=10$ to about
$2.8\times$ at $|C|=200$.

\begin{figure}[t]
  \centering
  \includegraphics[width=\linewidth]{figure2_runtime_ratio_mus_count}
  \caption{(a)~Greedy runtime relative to exact Z3. Dataset~B consistently
    exceeds $2\times$ owing to 5 sequential MUS-extraction passes.
    (b)~MUSes discovered by lazy greedy per instance. MUS count is constant
    across all sizes for each dataset, confirming that conflict complexity
    does not scale with $|C|$ in the current generators.}
  \label{fig:overhead}
\end{figure}

\subsection{Interpretation: Generator Simplicity, Not Algorithm Optimality}

The absence of an approximation gap is a property of the generators, not a
general property of the greedy algorithm.
All three datasets produce MUS hypergraphs that are structurally easy for
weighted greedy: Dataset~A has a unique two-element MUS with equal weights;
Dataset~B has a regular clique-like structure on one variable; Dataset~C
reduces to a two-element MUS after constraint propagation.

Adversarial instances for Set-Cover-style greedy require MUSes with disjoint
support and carefully constructed weights so that the greedy choice at each
step is locally cheap but globally suboptimal.
None of the three generator families produce such structure, because the
planted conflicts are confined to one or two variables with straightforward
weight assignments.
Constructing generators that force a gap---for example, by planting multiple
disjoint MUSes with cross-cutting weight gradients inspired by the Set Cover
hardness template of~\cite{feige1998}---is a natural direction for future work.
